% Page Layout
% Demonstrates: geometry, fancyhdr, two-column layout, margin notes.
% Compile with: pdflatex main.tex

\documentclass[12pt, a4paper]{article}
\usepackage[utf8]{inputenc}
\usepackage[T1]{fontenc}

% --- Page geometry: custom margins ---
\usepackage[
    top=2.5cm,
    bottom=2.5cm,
    left=3cm,
    right=3cm,
    marginparwidth=2cm,   % width for margin notes
]{geometry}

% --- Headers and footers ---
\usepackage{fancyhdr}
\pagestyle{fancy}
\fancyhf{}                              % Clear default headers/footers
\fancyhead[L]{\textit{LaTeX Tutorial}}  % Left header
\fancyhead[R]{\thepage}                 % Right header: page number
\fancyfoot[C]{\footnotesize Page Layout Example}  % Center footer
\renewcommand{\headrulewidth}{0.4pt}    % Header rule
\renewcommand{\footrulewidth}{0.4pt}    % Footer rule

% --- Packages for demo content ---
\usepackage{lipsum}        % Dummy text
\usepackage{marginnote}    % Better margin notes
\usepackage{multicol}      % Multi-column environments

\title{Page Layout and Typography}
\author{LaTeX Tutorial}
\date{\today}

\begin{document}
\maketitle
\thispagestyle{fancy}  % Apply fancy style to title page too

% =============================================================================
\section{Custom Margins with \texttt{geometry}}
% =============================================================================

The \texttt{geometry} package controls page dimensions. This document uses:

\begin{verbatim}
\usepackage[top=2.5cm, bottom=2.5cm,
            left=3cm, right=3cm]{geometry}
\end{verbatim}

\lipsum[1]

% =============================================================================
\section{Headers and Footers with \texttt{fancyhdr}}
% =============================================================================

Look at the top and bottom of this page. The header shows the document title
on the left and page number on the right. The footer has centered text.

\marginnote{\scriptsize This is a margin note! Use them for asides.}[0cm]

Key commands:
\begin{itemize}
    \item \verb|\fancyhead[L]{...}| --- left header
    \item \verb|\fancyhead[R]{...}| --- right header
    \item \verb|\fancyfoot[C]{...}| --- center footer
    \item \verb|\renewcommand{\headrulewidth}{0.4pt}| --- header line
\end{itemize}

\lipsum[2]

% =============================================================================
\section{Two-Column Layout}
% =============================================================================

You can switch to two columns for part of a document using the
\texttt{multicol} environment:

\begin{multicols}{2}
\lipsum[3]

\marginnote{\scriptsize Another margin note, visible in multi-column mode.}

\lipsum[4]
\end{multicols}

% =============================================================================
\section{Margin Notes}
% =============================================================================

Margin notes\marginnote{\scriptsize $E = mc^2$} are useful for annotations,
definitions, or quick references. Use \verb|\marginnote{text}| from the
\texttt{marginnote} package (more reliable than \verb|\marginpar|).

\lipsum[5]

% =============================================================================
\section{Three-Column Layout}
% =============================================================================

\begin{multicols}{3}[\subsection*{Three narrow columns}]
\small
\lipsum[6]
\end{multicols}

\end{document}
